\section{Difusión de la capacidad de cómputo: Capacity, Inference y brecha digital}

La combinación de modelos termina por apoyarse en la infraestructura. La clase calificó la inversión a gran escala como no-regret, pero desde el punto de vista sistémico ese juicio solo se sostiene si la capacity se aprovecha, la inference es asequible, las aplicaciones pueden entrar y los beneficios pueden retroalimentarse. Esta sección usa el flywheel para analizar cómo tiene éxito «comprar capacidad de cómputo» y explica también por qué puede ampliar, en lugar de reducir, la brecha digital.

\subsection{Utilization flywheel y token economics}

La sección anterior resolvió «qué modelo elegir»; esta responde «cómo lograr que la capacity a escala nacional se use de verdad». Al leer la figura hay que distinguir entre installed capacity, available service, utilization real y application demand: solo cuando mejoran a la vez la planificación, el onboarding de modelos, la confiabilidad y la adopción por parte de los usuarios puede bajar el unit cost y, a su vez, atraer más workload efectivo; añadir racks por sí solo no pone en marcha el volante.

\lecturefigure{11-infrastructure-flywheel.png}{La capacity solo forma una retroalimentación positiva a través de la utilization, el unit cost y la application demand.}{Subtítulos locales 27:10--32:20, 42:30--45:10; contexto oficial de Groq/Aramco Digital.}

\begin{knowledgebox}{Lectura de la figura: el token price más bajo no es el único objetivo}
Un precio bajo puede ampliar la experimentation, pero los usuarios también necesitan model availability, calidad en árabe y en dominios verticales, latency, data protection, reliability y support. Si el workload no es estable o el stack de software es difícil de usar, la capacity se convierte en un subsidio; si la demanda de aplicaciones crece y la power y la network no tienen margen, la promesa de precios bajos puede a su vez perjudicar el SLO.
\end{knowledgebox}

\begin{warningbox}{Hay que separar el momento de la clase de los anuncios posteriores}
La entrevista trata la colaboración de inferencia entre Groq y Aramco Digital tal como era visible en marzo de 2025; el anuncio posterior de ampliación en LEAP 2025 puede servir para ilustrar la evolución, pero no debe presentarse retroactivamente como capacidad que ya se había entregado en el momento de la clase. Por eso estas notas solo toman como conclusión del mecanismo «incorporar rápidamente múltiples modelos mediante onboarding y reducir la barrera de entrada a la inference».
\end{warningbox}

\subsection{Training e inference son dos formas de infraestructura}

El \term{training} (entrenamiento) suele actualizar parámetros mediante trabajos síncronos a gran escala y busca rendimiento y time-to-train; la \term{inference} (inferencia) usa un modelo ya entrenado para responder solicitudes y está restringida por latency, availability, procesamiento por lotes, KV cache, ubicación geográfica y costos continuos. Ambos comparten accelerators, pero no comparten exactamente la misma topology, el mismo scheduler ni la misma business demand.

\lecturefigure{12-training-inference.png}{El training es más concentrado y síncrono; la inference es más distribuida y está restringida por el SLO.}{Subtítulos locales 32:20--36:50; redibujo conceptual.}

\begin{knowledgebox}{Lectura de la figura: no fijar una proporción permanente}
El juicio de la clase sobre la proporción training/inference es una observación de un momento dado. El desarrollo de modelos, el fine-tuning, los synthetic data, el test-time compute, los agent loops y el crecimiento de usuarios cambian el workload mix. Un capacity plan debe usar un scenario range y conjuntos modulares, en lugar de extrapolar la proporción actual hasta el final de la vida útil de las instalaciones.
\end{knowledgebox}

\begin{importantbox}{La versión de recursos y la versión de capacidades de la brecha digital}
La brecha en su versión de recursos consiste en el access price, la network y la accelerator availability; la brecha en su versión de capacidades consiste en datos, talento, workflow redesign, evaluation y procurement. Reducir solo el token price puede mejorar la primera, pero no mejora automáticamente la segunda. Los programas públicos deben incluir el adoption support y el outcome measurement en el presupuesto de infraestructura.
\end{importantbox}

\subsection{Resumen de la sección}

La inversión en capacidad de cómputo solo se convierte en un productive asset cuando entra en el utilization flywheel. El training y la inference tienen formas de recursos distintas, y la inclusión digital no se reduce a tokens baratos; el software, las skills, los datos y la organizational adoption también forman parte de la infraestructura.
